% Chapter 1: Introduction (Massive Expansion for 80+ Page Project)
\chapter{INTRODUCTION}
%% \thispagestyle{plain} removed to allow global fancy style with logo

The Cloud First era has revolutionized the digital economy in the world. With organizations progressing beyond mere cloud adoption to Multi-Cloud Strategic Maturity, the necessity to integrate management platforms has emerged as a serious operational necessity.

\section{Historical Context of Computing}
There have been a number of paradigms of transformation in modern computing, all of which have made the user and the underlying hardware more and more abstract. It is important to understand how these paradigms have developed to fully recognize the value proposition of the Nebula Multi-Cloud Command Center.

\subsection{The Mainframe Era (1950s - 1970s)}
Early computer technology saw the use of massive computers which occupied whole rooms. These systems were centralized and users accessed these systems through dumb terminals. Though very powerful at the time, they were prohibitively priced, difficult to scale, and were a single point of failure. It was highly manualized management that demanded a specialized workforce.

\subsection{Client-Server Architecture (1980s - 1990s)}
The client-server model was the result of the emergence of the Personal Computer (PC). Calculations were performed over a server (where data or services are delivered) and a client (where they are requested). Although more flexible compared to the mainframes, this period brought about the so-called silo issue where various departments operated with their own physical servers, resulting in inefficient use of the hardware resources and excessive maintenance expenses.

\subsection{The Virtualization Revolution (2000s)}
The relationship between software and hardware was changed by virtualization. Hyper- visors enabled more than one Virtual Machines (VMs) to operate on a physical server. It was the technological antecedent to cloud computing and allowed efficient resource utilization and the ability to spin up servers in minutes as opposed to weeks.

\section{The Rise of Cloud Computing}
Cloud computing is the natural extension of virtualization, wide network access and automated management. According to the National Institute of Standards and Technology (NIST) [1], cloud computing refers to a model to facilitate ubiquitous, convenient, on-demand network access to a common pool of configurable computing resources.

\subsection{NIST Service Models}
To get the picture of what Nebula platform covers, we need to classify the services that it deals with:
\begin{enumerate}
    \item \textbf{Infrastructure as Service (IaaS)}: Offers basic computing, including Virtual Machines (AWS EC2, Azure VM), storage and networks. This layer is heavily integrated by Nebula.

    \item \textbf{Platform as Service (PaaS)}: It is a platform that enables customers to create, execute, and manage applications without having the complexity of infrastructure. Examples are AWS Lambda and Azure Functions, which Nebula is managing as Network Functions.
    \item \textbf{Software as Service (SaaS)}: Software is sold on a subscription model and is centrally hosted. Although Nebula is a SaaS-like platform, it controls the underlying IaaS and PaaS to facilitate enterprise SaaS deployments.
\end{enumerate}

\subsection{NIST Deployment Models}
\begin{itemize}
    \item \textbf{Public Cloud}: It is owned and managed by a third party service provider (AWS, Azure, GCP).
    \item \textbf{Private Cloud}: Resources known to only one business or organization.
    \item \textbf{Hybrid Cloud}: A combination of two or more clouds (community, private, or public) that are distinct entities, yet are interconnected.
    \item \textbf{Multi-Cloud}: This is where a number of public clouds of various providers are used. Nebula is actually created to conquer this space.
\end{itemize}

\section{The Multi-Cloud Strategic Shift}
A multi-cloud strategy is another approach that organizations are employing to prevent Vendor Lock- in. Nonetheless, this transition presents important managerial strife. 

\subsection{Drivers for Multi-Cloud Adoption}
\begin{itemize}
    \item \textbf{Best-of-Breed Services}: Data analytics with GCP and retaining enter- prise apps on Azure.
    \item \textbf{Data Sovereignty}:Storing data in particular physical locations to abide by GDPR or local laws.
    \item \textbf{Negotiation Leverage}: Presence in several clouds eliminates dependency on one vendor and their pricing model.
\end{itemize}

\section{Problem Statement and Motivation}
Multi-cloud management is disjointed, despite its advantages. At this time, an administrator should:
\begin{enumerate}
    \item Control high-security credentials in three portals.
    \item Get familiar with different naming (e.g., Security Groups in AWS and NSG in Azure).
    \item Manage deployments that are long-running and frequently crash in the browser.

\end{enumerate}
Nebula is motivated by the need to develop a Unified Control Plane which encapsulates these complexities into one, high-performance API and User Interface.

\section{Project Objectives and Scope}
The particular targets of this project are:
\begin{itemize}
    \item Achieve 100% visibility of cloud inventory across providers.
    \item Minimize troubleshooting with LLM log analysis.
    \item Automatize resource provisioning processes.
    \item A secure credential vault by way of application-layer encryption.
    \item Stream the deployment logs in real time with WebSockets or polling.
\end{itemize}
The area includes the early work on the core engine, security vault and a React-based interactive cockpit.

\newpage
